\subsection{Motivierendes Beispiel}
\begin{frame}[fragile]{Prädikatenlogik among us}
	\begin{itemize}
		\item<1-> Ihr seid auf einem Raumschiff
		\item<2-> Mit eurer Crew
		      $$
			      U = \{\sus{red}, \sus{green}, \sus{blue}, \sus{white}, \sus{yellow}, \sus{magenta}\}
		      $$
		\item<3-> Aber manche von euch sind Imposter

	\end{itemize}
	\only<2>{\hspace{3cm}\includegraphics[width=4.5cm]{../figures/amongus_crew.png}}
	\only<3>{\hspace{3cm}\includegraphics[width=4.5cm]{../figures/amongus_imposter.jpg}}
\end{frame}

\begin{frame}[fragile]{Prädikatenlogik among us}
	\begin{itemize}
		\item<1-> Wir brauchen eine Funktion $C$ um zu erkennen wer ein sicheres Crewmitglied zu erkennen
		      $$C: U \to \{wahr, falsch\}$$
		\item<2-> Und eine Funktion $A$ um zu beschreiben ob ein Crewmitglied noch lebt
		      $$A: U \to \{wahr, falsch\}$$
		\item<3-> Und eine Funktion $M$ um zu beschreiben ob zwei Spieler aufeinander treffen
		      $$M: U\times U \to \{wahr, falsch\}$$
	\end{itemize}
\end{frame}


\begin{frame}[fragile]{Eine kurze Geschichte}
	$$
		\overline{C(\sus{red})}\wedge C(\sus{white})\wedge M(\sus{red},\sus{white})\implies \overline{A(\sus{white})}
	$$
\end{frame}

\begin{frame}[fragile]{Aber wer ist der Imposter?}
	\begin{itemize}
		\item<1-> \alert{Problem:} Wer ist der Imposter?
		\item<2-> Wir wissen nicht welche Crewmitglieder wir für was einsetzen können
		\item<3-> Also müssen wir allgemeiner werden:
		      $$
			      \exists x,z \in U: \overline{C(x)}\wedge C(y)\wedge M(x,y)\implies \overline{A(y)}
		      $$
	\end{itemize}
	\begin{center}
		\only<4>{Das ist eine prädikatenlogische Formel}
	\end{center}
\end{frame}

\subsection{Definitionen}

\begin{frame}{Prädikatenlogik erster Ordnung}
	\begin{itemize}
		\item<1-> \alert{Variablen} sind wie ihr sie aus der Aussagenlogik bereits kennt.
		\item<2-> $U$ ist unser \alert{Universum}. Es enthält alle Elemente die wir in unsere Variablen einsetzen dürfen
		\item<3-> \alert{Prädikatensymbole} sind Mappings von Elementen aus $U$ auf $wahr$ oder $falsch$.
		      Sie werden normal groß geschrieben.
		\item<4-> \alert{Funktionssymbole} sind Mappings von Elementen aus $U$ auf andere Elemente aus $U$.
		      Sie werden normal klein geschrieben.
		\item<5> Eine \alert{Interpretation} $I$ ist ein Lösungsvorschlag für die prädikatenlogische Formel.
		      Es ist ein Mapping von Variablen auf Elemente des Universums, der Prädikatensymbole auf Prädikate und der Funktionsformeln auf Funktionen.
	\end{itemize}
\end{frame}

\begin{frame}{Unser Beispiel}
	\begin{itemize}
		\item $U = \{\sus{red}, \sus{green}, \sus{blue}, \sus{white}, \sus{yellow}, \sus{magenta}\}$
		\item Prädikate:
		      \begin{itemize}
			      \item $P_{sus}(x,y)$: $x$ verdächtigt $y$
			      \item $P_{kill}(x,y)$: $x$ tötet $y$
			      \item $P_{seen}(x,y)$: $x$ und $y$ wurden zusammen gesehen
			      \item $P_{alive}(x)$: $x$ lebt
			      \item $P_{crew}(x)$: $x$ ist ein echtes Crewmitglied
		      \end{itemize}
	\end{itemize}
\end{frame}

\begin{frame}{Beispiel}
	$$P_1(x,y) \implies \overline{P_2(x)}$$
	\pause
	\alert<only:+->{Interpretation:}
	\begin{align*}
		I(P_1) & = P_{kills}   \\
		I(f_2) & = P_{crew}    \\
		I(x)   & = \sus{red}   \\
		I(y)   & = \sus{white}
	\end{align*}
	\pause
	\sus{red} tötet \sus{white}. Dann muss \sus{red} ein Imposter sein!
\end{frame}

{\setbeamercolor{palette primary}{bg=ExColor}
\begin{frame}[fragile]{Denkpause}
	\footnotesize
	\begin{alertblock}{Aufgaben}
		Findet Interpretationen für die folgenden Prädikatenlogischen Formeln in unserem Beispiel:
	\end{alertblock}
	\metroset{block=fill}
	\begin{block}{Normal}
		\begin{itemize}
			\item $P_1(x,y) \Leftrightarrow P_2(y,x)$
			\item $(P_1(x,y) \wedge \overline{P_2(y)}) \implies P_3(y)$
		\end{itemize}
	\end{block}
	\begin{block}{Etwas Schwerer}
		\begin{itemize}
			\item $(P_1(x,y) \wedge P_2(y)) \implies P_3(x)$
		\end{itemize}
	\end{block}
	\alert{Hinweis:} In dem Spiel Among Us gibt es viele Gründe etwas zu tun oder nicht zu tun.
	Verschwendet nicht zu viel Zeit die Aktionen zu überdenken. Wenn ihr argumentieren könnt warum jemand agiert wird es schon passen.
\end{frame}

\begin{frame}{Lösung}
	$$(P_1(x,y) \wedge \overline{P_2(y)}) \implies P_3(y)$$
	\begin{align*}
		I(P_1) & = P_{seen}     \\
		I(P_2) & = P_{seen}     \\
		I(x)   & = \sus{blue}   \\
		I(y)   & = \sus{yellow}
	\end{align*}
	Wir haben \sus{red} und \sus{green} zusammen gesehen und \sus{green} ist tot. Dann muss \sus{red} ein Imposter sein.
\end{frame}

\begin{frame}{Lösung}
	$$(P_1(x,y) \wedge \overline{P_2(y)}) \implies P_3(y)$$
	\begin{align*}
		I(P_1) & = P_{seen}    \\
		I(P_2) & = P_{alive}   \\
		I(P_3) & = P_{crew}    \\
		I(x)   & = \sus{red}   \\
		I(y)   & = \sus{green}
	\end{align*}
	Wir haben \sus{red} und \sus{green} zusammen gesehen und \sus{green} ist tot. Dann muss \sus{red} ein Imposter sein.
\end{frame}

\begin{frame}{Lösung}
	$$(P_1(x,y) \wedge P_2(y)) \implies P_3(x)$$
	\begin{align*}
		I(P_1) & = P_{sus}       \\
		I(P_2) & = P_{alive}     \\
		I(P_3) & = P_{crew}      \\
		I(x)   & = \sus{white}   \\
		I(y)   & = \sus{magenta}
	\end{align*}
	\sus{white} verdächtigt \sus{magenta}, und \sus{white} ist noch am Leben. Also ist \sus{magenta} kein Imposter.
\end{frame}
}



\begin{frame}[fragile]{Reminder: Quantoren}
	Mithilfe von \textbf{Quantoren} vereinfachen wir uns die Schreibweise dieser Aussagen.\\
	\vspace{0.5cm}
	Quantor \alert{$\forall$}: Die Aussage gilt für alle Elemente.\\
	\metroset{block=fill}
	\begin{exampleblock}{Beispiel}
		$A_1$: $\forall k \in \mathbb{N}:$ $2k$ ist gerade
	\end{exampleblock}
	Quantor \alert{$\exists$}: Die Aussage gilt für mindestens ein Element.\\
	\metroset{block=fill}
	\begin{exampleblock}{Beispiel}
		$A_2$: $\exists k \in \mathbb{N}:$ $k$ ist Primzahl
	\end{exampleblock}
\end{frame}

\begin{frame}{Funktionen}
	\begin{itemize}
		\item Funktionen erlauben die \alert{Transformation von Elementen} zu beschreiben
		\item<2-> Sie wandelnt Elemente aus $U$ in ein Element aus $U$ um.
		\item<3-> Funktionssymbole werden normal durch Kleinbuchstaben representiert
	\end{itemize}

\end{frame}

\begin{frame}{Ein weiteres Beispiel}
	\begin{itemize}
		\item $U = \mathbb Z$
		\item Prädikate:
		      \begin{itemize}
			      \item $P_{ood}(x)$: $x$ ist ungerade
			      \item $P_{prim}(x)$: $x$ ist eine Primzahl
			      \item $P_{+}(x)$: $x$ ist positiv
			      \item $P_>(x,y): x>y$
			      \item $P_=(x,y): x=y$
		      \end{itemize}
		\item Funktionen: $+, -, \cdot$
	\end{itemize}
\end{frame}

\begin{frame}{Ein weiteres Beispiel}
	$$\exists x: P_1(f_1(x,x),x)$$
	\pause
	\begin{align*}
		I(P_1) & = P_{=} \\
		I(f_1) & = P_{+} \\
	\end{align*}
	\pause
	Es existiert ein $x \in \mathbb N$ so dass $x\cdot x = x$
	(Das wäre dann die $1$ oder die $0$)
\end{frame}

{\setbeamercolor{palette primary}{bg=ExColor}
\begin{frame}[fragile]{Denkpause}
	\footnotesize
	\begin{alertblock}{Aufgaben}
		Findet Interpreatationen für die folgenden Prädikatenlogischen Formeln in unserem Beispiel:
	\end{alertblock}
	\metroset{block=fill}
	\begin{block}{Normal}
		\begin{itemize}
			\item $\forall x,y : P_1(x) \wedge P_1(y) \implies \overline{P_1(f_1(x,y))}$
			\item $\forall x\exists y: P_1(f_1(x,y),x)$
			\item $\forall x\exists y: P_1(f_1(x,y),y)$
		\end{itemize}
	\end{block}
	\begin{block}{Etwas Schwerer}
		\begin{itemize}
			\item $\exists x,y: P_1(x,y) \wedge P_2(f_1(x,y),x) \wedge \overline{P_2(f_1(y,x),x)}$
		\end{itemize}
	\end{block}
\end{frame}

\begin{frame}{Lösung}
	$$\forall x,y : P_1(x) \wedge P_1(y) \implies \overline{P_1(f_1(x,y))}$$
	\begin{align*}
		I(P_1) & = P_{ood} \\
		I(f_1) & = +
	\end{align*}
	Für alle $x,y \in \mathbb Z$ gilt: wenn $x$ und $x$ ungerade sind, dann ist $x+y$ ungerade
\end{frame}

\begin{frame}{Lösung}
	$$\forall x\exists y: P_1(f_1(x,y),x)$$
	\begin{align*}
		I(P_1) & = P_{=} \\
		I(f_1) & = +
	\end{align*}
	Für alle $x \in \mathbb Z$ existiert ein $y$ so dass $x+y = x$. (Dieses $y$ ist $0$).
\end{frame}

\begin{frame}{Lösung}
	$$\forall x\exists y: P_1(f_1(x,y),y)$$
	\begin{align*}
		I(P_1) & = P_{=} \\
		I(f_1) & = \cdot
	\end{align*}
	Für alle $x \in \mathbb Z$ existiert ein $y$ so dass $x\cdot y = y$. (Dieses $y$ ist $0$).
\end{frame}

\begin{frame}{Lösung}
	$$\exists x,y: P_1(x,y) \wedge P_2(f_1(x,y),x) \wedge \overline{P_2(f_1(y,x),x)}$$

	\begin{align*}
		I(P_1) & = P_{>} \\
		I(P_2) & = P_{=} \\
		I(f_1) & = -
	\end{align*}

	Es exsitieren zwei Zahlen $x,y \in \mathbb Z$ so dass $x$ größer ist als $y$ und $x-y=x$ aber $y-x\neq x$

\end{frame}
}



